\documentclass{article}
\usepackage{/globals/de}
\usepackage{/globals/laws}
\usepackage{/globals/anki}
\ankiprefix{bundesrat}
\ankitopic{Bundesrat}
\ankishow
\begin{document}
\section{Der Bundesrat} 
Der Bundesrat gibt den Bundesländern die möglichkeit an Gesetzen \uae mitzuwirken.
\ankinote{sinn}{Sinn}{Länder mitwirkung geben}

Dabei hat er im Gesetzgebungsprozess aber eine geringere Bedeutung.

Der Bundesrat ist ein ewiges Organ weil es nicht wie der Bundestag Wahlperioden oder ein Diskontinuitätsprinzip hat.
\ankinote{ewigkeit}{Ewigkeit}{Eweig, keine Wahlperioden, kein Diskontinuitätsprinzip}

Als Rechtsquelle nutzt der Bundesrat \art{23} \GG, \arts{50}{53} \GG, \gesetz{EUZBLG} und die \GOBR.
\ankinote{rechtsq}{Rechtsquellen}{Art. 23 GG, Art 50-53 GG, EUZBLG, GO BR}

\subsection{Aufgaben}
Zu den Aufgaben gehören die Mitwirkung an der Gesetzesgebung von Bundesgesetzen, der Verwaltung des Bundes und bei manchen EU-Angelegenheiten.
\ankinote{aufgaben}{Aufgaben}{Mitwirkung Bundesgesetze, Verwaltung, EU-Angelegenheit}

Zu den EU-Angelegenheiten zählen EU-Gesetze die Aspekte betreffen die eigentlich auf Länderebene gehandhabt werden, wie das Schulwesen.

\subsection{Mitwirkung bei Gesetzen}
Der Bundesrat hat ein Initiativrecht für Bundesgesetze und grobere Forderungen durch sogenannte \emph{Entschließungen}.

Gesetzesentwürfen kann der Bundesrat je nach Gesetzesart zustimmen oder Einspruch einlegen, oder einen Vermittlungsausschuss ausrufen.

Der Bundesrat kann Stellung nehmen.
\ankinote{mitwirkung-g}{Mitwirkung Gesetze}{Initiativrecht, Entschließungen, Stellung, Zustimmen/Einspruch}

\subsection{Mitwirkung bei der Verwaltung}
Der Bundesrat kan Verordnungen \uae zustimmen und dabei \sog \emph{Maßgaben} verlangen; bestimmte Veränderungen an der Verordnung.
\ankinote{mitwirkung-v}{Mitwirkung Verwaltung}{Zustimmen unter Maßgaben (Bedingungen)}

\subsection{Zusammensetzung}
Der Bundesrat hat insgesamt 69 Stimmen. Dabei haben größere Bundesländer mehr Stimmen, aber nicht proportional; für die föderale Gleichheit.
\ankinote{sitze}{Sitze}{69, nach Bundesland größe, föderale Gleichheit}

Die MitgliederInnen sind von den Ländern gewählte MitgliederInnen derer Landesregierungen. Sie haben im Gegensatz zu MdBs kein freies Mandat, sondern ein \emph{imperatives} Mandat; alle Stimmen eines Bundeslandes müssen einheitlich sein. Dafür wird oftmals ein \sog \emph{Stimmführer} genutzt der wählt, so wie es die Landesregierung festgelegt hat.
\ankinote{mitglieder}{Mitglieder}{Aus Landesregierungen gesendet}
\ankinote{mandat}{Mandat}{imperativ, Stimmführer wählt}

Koalitionen entscheiden sich \idr entweder dazu sich bei Uneinigkeiten zu enthalten oder nur bei eher relevanteren Uneinigkeiten zu enthalten.

\subsection{Arbeitsweise}
Der Bundesrat hat seltene Sitzungen in einem öffentlichen Plenum, hat nicht öffentliche Ausschüsse, ein Präsidium mit einem ständigen Rat und eine Europakammer.
\ankinote{arbeitsw}{Arbeitsweise}{Seltesnes öffentliches Plenum, privae Ausschüsse, Präsidium, ständiger Rat, Europakammer}
\end{document}